\begingroup
    % Suppress page breaks for chapter commands within this group
    \let\clearpage\relax
    \let\cleardoublepage\relax
    
    % Acknowledgements Section
    \phantomsection
    \addcontentsline{toc}{chapter}{Acknowledgements}
    \chapter*{\centering{\Large Acknowledgements}}
    \vspace{-2em}
    I would like to thank my supervisor, Prof. Pulin Gong, for offering such an engrossing research project, balancing exciting results with an emphasis on the \textit{so-what} of it all, and providing guidance across every area of research in my Honours year and in my undergraduate studies. I would also like to thank my co-supervisor, Dr. Brendan Harris, for the endless help and guidance in the most important areas of research, and for introducing me to the Julia hype. I am thankful of my broader research group, for bringing the Madsen Honours office closer to Physics. Among the Physics Honours cohort, I'm particularly glad to be a part of a smaller number who almost always find themselves at the SiFA research office, or else on the short walk to Ralph's café. Finally, and most importantly, I cannot thank enough my family and friends, across Australia and Germany, who---no matter how hard I've made it---always give their unceasing support at any distance.
    
    \vspace{3em} % spacing between sections
    
    \phantomsection\addcontentsline{toc}{chapter}{Statement of student contribution} % Don't remove me!
    \chapter*{\centering{\Large Statement of student contribution}}
    \vspace{-2em}
    The central idea of this project was devised by my supervisor, Prof. Pulin Gong, and myself. The theoretical framework of fractional neural sampling (FNS) was previously developed by my supervisor to characterise the rich spatiotemporal dynamics of cortical neural activity. Building upon this theoretical foundation, my supervisor proposed extending the FNS framework to model gaze dynamics. My supervisor also provided a deep understanding of the literature involved---spanning the fields of physics, neuroscience, and machine learning---as well as providing the collaboration with Dr. Pascal Fries, Tim Näher, and Yufeng Zhang, who were invaluable in providing the necessary data and experimental framework to evaluate this work. My co-supervisor Dr. Brendan Harris provided a rich understanding of FNS theory as concerns cortical activity. My co-supervisor also provided practical advice as to computational implementation and tools to be used. My co-supervior provided the basis for the code which structures the numerical solving of the FNS equation in the Julia programming language.
    
    My original contributions to this work are as follows. Broadly, my conceptual additions to this project were in the addition of an informative kernel to couple information gathering to model dynamics, and the interpretation of a framework to span spatiotemporal scales of gaze. Additionally, considering the explored dynamical landscape as a dynamical Bayesian belief distribution was also introduced by myself. The rediscovery of OMT in the eye-tracking data and the model was initiated by myself. Specifically, I provided the careful analysis of kernel-driven rhythmicity, local harmonic oscillations, and Lévy superdiffusion in this system (\cref{chap:a_unifying_model}). The truncated solution of the fractional Fokker-Planck equation in this system, and dynamical approximations to find the theta- and gamma-timescales was done by myself. The writing of the scripts to simulate and analyse the modelled system in the Julia programming language was also done by myself, in the \texttt{VisualFNS.jl} project. I devised the methodology to compare the model with eye-tracking data (\cref{chap:empirical_validations}). This includes the $q$--order fractional moments compared between the model and the data, which I interpreted as multiscale superdiffusion, through research in literature on SAND. This includes the saccade detection from modelled solutions, as presented and documented in the \texttt{VisualFNS\_Analysis.jl} project. The construction and methodology of constructing informational metrics was done by myself (\cref{chap:functional_advantages}). I defined the quantification of efficiency, under and guidance of my supervisors, and provided the interpretation of each considered dynamical signature and compared each to generalised forms of probabilistic computation. The code for defining and calculating these metrics can be found in the \texttt{VisualData\_Analysis.jl} project.

    Generative AI was not used to construct original code, develop theory, or write up and analyse results. Generative AI was used to debug code and in the initial stages of literature search.

\endgroup
\newpage % Force a new page after both sections are finished